
\section{Introduction}
\label{sec:introduction}

\subsection{The Simulation Gap in Silicon Spin Qubit Control}
\label{subsec:simulation_gap}

Silicon spin qubits—particularly SiMOS devices—are poised to scale quantum
computing. With coherence times approaching $T_1 \sim 1$\,ms and
$T_2^* \sim 0.4$\,ms~\cite{Xue2022,Noiri2022,Philips2022}, CMOS
fabrication compatibility, and natural nearest-neighbour connectivity, they
offer a manufacturable path toward fault-tolerant architectures. Yet their
noise sensitivity remains a critical roadblock: charge fluctuations, hyperfine
coupling variations, and residual spin-orbit interactions degrade gate fidelity
in ways that first-principles prediction cannot reliably
capture~\cite{Veldhorst2015,Fogarty2018}.

Closing this gap demands simulation environments that are simultaneously
physically rigorous and accessible to the full spectrum of quantum control
algorithms. This matters for three concrete reasons. First, comparing control
paradigms on identical physical models is the only way to isolate algorithmic
performance from modelling artefacts—a prerequisite for principled algorithm
selection in hardware co-design. Second, a unified interface enables hybrid
strategies (e.g., gradient-initialised RL or BO-tuned GRAPE) that are otherwise
infeasible. Third, reproducibility requires that benchmark results be obtainable
by any researcher regardless of their preferred algorithm library.

No existing open-source platform satisfies all three requirements for silicon
spin qubits. Lindblad solvers such as QuTiP~\cite{Johansson2012,Johansson2013} deliver precise
open-system dynamics but expose no standard algorithm interface. Gradient-based
frameworks such as C3~\cite{Wittler2021} support GRAPE, CRAB, and CMA-ES but
exclude reinforcement learning and silicon-specific noise models. Quantum RL
environments~\cite{VanDerLinde2023,QiskitGym2024} offer Gymnasium compatibility
but rely on abstract depolarising noise rather than device-specific physics.
Qiskit Pulse~\cite{Alexander2020} enables pulse-level control for
superconducting qubits but omits SiMOS models entirely. TensorFlow
Quantum~\cite{Broughton2020} targets gate-model circuits via Cirq, not spin
qubit physics. Orbell's Oxford work~\cite{Orbell2024} advances Bayesian device
tuning for SiMOS but operates at device characterisation rather than gate
synthesis. No single integrated platform supports all five paradigms—gradient
(GRAPE, CRAB), variational (VQE), Bayesian (BO, CMA-ES), and deep RL—on the
same device-accurate silicon spin qubit model. SiliQun fills this gap as the
first open-source platform to unify device-accurate Lindblad simulation with a
standard algorithm-agnostic interface for silicon spin qubits.

\subsection{The Physics-Grounded Iterative Refinement Stack}
\label{subsec:intro_pgirs}

We introduce the \textbf{Physics-Grounded Iterative Refinement Stack} (PGIRS)
methodology for building simulation platforms that are both physically rigorous
and algorithm-agnostic. PGIRS rests on five design principles:

\begin{itemize}
  \item \textbf{Noise-First Design:} Noise models are calibrated from
    experimental device data and validated against analytical benchmarks before
    any control algorithm is integrated.
  \item \textbf{Standard Interface Compliance:} A Gymnasium-compatible
    \texttt{Env} interface provides plug-and-play access for any algorithm
    library—RL, gradient-based, or variational—without modification.
  \item \textbf{Reward Engineering:} A single reward function whose weights are
    optimised offline via Bayesian methods is fixed identically across all
    algorithms, ensuring comparability without conflating reward design with
    policy learning.
  \item \textbf{Hierarchical Decomposition:} Multi-qubit objectives decompose
    into sub-goals with dense intermediate rewards, enabling tractable
    optimisation for any algorithm.
  \item \textbf{Empirical Validation:} Each component undergoes analytical
    benchmarking, cross-algorithm testing, and hardware transfer checks before
    publication.
\end{itemize}

SiliQun implements PGIRS across the silicon spin qubit family. The key
distinction from prior work is that PGIRS treats the control algorithm as an
interchangeable component: the same physics engine, reward function, and
hardware profile serve gradient-based optimisers, variational circuits,
Bayesian search, and deep RL agents without modification.

\subsection{Contributions}
\label{subsec:contributions}

This work makes four primary contributions centred on the SiliQun simulation
platform, and one secondary contribution demonstrating its use:

\begin{enumerate}
  \item \textbf{SiliQun physics engine:} A unified Lindblad solver supporting
    SiMOS, donor-electron, GAA, Si/SiGe quantum dot, and SLEDGE qubits via
    configurable hardware profiles, validated at numerical precision
    ($\max\,\delta \leq 2.78 \times 10^{-16}$) and GPU-accelerated for
    $93\times$ faster simulation cycles.
  \item \textbf{Multi-technology hardware profiles:} A \texttt{TechnologyProfile}
    abstraction encoding device physics for five silicon-compatible spin qubit
    technologies, grounded in published experimental data and selectable via a
    single constructor argument.
  \item \textbf{Algorithm-agnostic control interface:} A Gymnasium-compatible
    environment with a discrete 11-token gate action space, 16-element
    density-matrix state representation, and a pre-validated reward function
    applied identically to any control algorithm without modification.
  \item \textbf{Qiskit BackendV2 provider:} The pip-installable
    \texttt{siliqun-provider} exposes SiliQun as a Qiskit
    \texttt{BackendV2}-compatible backend, integrating directly with the Qiskit
    ecosystem.
  \item \textbf{Cross-paradigm benchmark:} We benchmark representative
    algorithms from five quantum control paradigms—gradient-based (GRAPE, CRAB),
    variational (VQE), Bayesian (BO, CMA-ES), and deep RL (DQN, SAC)—on
    identical hardware profiles, providing the first direct cross-paradigm
    comparison on a device-accurate silicon spin qubit simulator.
\end{enumerate}
